\section*{Chapter \ref{ch:ll:linux-tuning} --- Linux Tuning}

\begin{solution}{exo:ll:linux-tuning:1}
$1\,000 \times \qty{2}{\micro\second} = 0.2\%$ of the time. The mean added wait is $1\,000 \times (\qty{2}{\micro\second})^2 /
(2 \times \qty{1}{\second}) = \qty{2}{\nano\second}$. On a tickless CPU running one task, the tick no longer interrupts the
thread: the worst case loses the tick's full duration, \qty{2}{\micro\second} here.
\end{solution}

\begin{solution}{exo:ll:linux-tuning:2}
The call fails with \texttt{ENOMEM}: an ordinary user's locked-memory limit is 64~MiB on this laptop, and 3~GiB is far above
it. Raise the process's \texttt{RLIMIT\_MEMLOCK} (in the service's configuration or the system's limits), or give it the
\texttt{CAP\_IPC\_LOCK} capability, for which the limit is not enforced; the audit's \texttt{memlock} rule checks the first.
\end{solution}

\begin{solution}{exo:ll:linux-tuning:3}
\texttt{isolcpus=12-15,44-47} and \texttt{nohz\_full=12-15} (with \texttt{rcu\_nocbs=12-15}). The siblings 44--47 share the
hot cores' execution units: take them offline, or leave them isolated with nothing pinned to them, so that no other thread
competes with the hot threads for their cores.
\end{solution}

\begin{solution}{exo:ll:linux-tuning:4}
Both steal 1\%. The thousand short gaps add a mean wait of $1\,000 \times (10^{-5})^2/2 = \qty{50}{\nano\second}$, and 0.5\%
of messages wait more than \qty{5}{\micro\second}. The single \qty{10}{\milli\second} gap adds $(10^{-2})^2/2$ seconds,
\qty{50}{\micro\second} on average, a thousand times more, and almost 1\% of messages wait more than
\qty{5}{\micro\second}, some of them ten milliseconds.
\end{solution}

\begin{solution}{exo:ll:linux-tuning:5}
When the microseconds saved are worth more than a core, its power and its heat: on the few threads of the hot path, not on
the rest. Besides the core, it costs power and heat (the kernel documentation warns that polling in the idle loop can overheat
the processor and disables turbo frequencies), and on a hyper-threaded core it slows whatever runs on the sibling.
\end{solution}

\begin{solution}{exo:ll:linux-tuning:6}
With default throttling, real-time threads may use \qty{950}{\milli\second} of each second; when the agent is runnable, it gets
the other \qty{50}{\milli\second}: a gap of up to \qty{50}{\milli\second} once a second, 5\% of the time. Fixes: move the agent
off the CPU (isolate it), after which the throttle has nothing to give the time to; or disable throttling, which risks locking
up the CPU if the thread misbehaves. Better, run the feed thread as an ordinary thread on an isolated core.
\end{solution}

\begin{solution}{exo:ll:linux-tuning:7}
Parse the fixture's \texttt{/proc/interrupts} for lines whose name contains the interface (for example \texttt{ens1f0-TxRx-0}),
and require each such line's \texttt{smp\_affinity\_list} to be a subset of the housekeeping CPUs of the plan's node. The tuned
fixture binds them to CPUs 16--20 and passes; add a line to the mistuned fixture bound to CPU 3 (node 0) and assert the
failure names it.
\end{solution}

\begin{solution}{exo:ll:linux-tuning:8}
Isolated CPUs are taken out of the scheduler's load balancing: threads whose affinity covers several isolated CPUs are not
spread over them and can all stay on the one where they started, time-sharing it (the shared curve of
\cref{fig:ll:linux-tuning:hiccups}). Pin each thread to its own isolated CPU, as \texttt{firm.affinity} does from the plan.
\end{solution}

\begin{solution}{pb:ll:linux-tuning:1}
\begin{enumerate}
\item Either would bring the kernel onto the CPU (a \omterm{def:ll:cpp-for-latency-i-memory:fault}{page fault}, a system call), and the meter would measure itself instead of
what is imposed on it.
\item On the committed run, 5 seconds for about 323 million passes: about \qty{15}{\nano\second} a pass, so a
\qty{200}{\nano\second} gap is more than ten passes lost, not a slow iteration.
\item About 1.3\%, with a worst gap of about \qty{5.4}{\milli\second}.
\item About 51\%: two equal threads that both always want the CPU each get half of it.
\item About \qty{3.1}{\micro\second} on the committed run; it varies with the run's longest gap.
\item About 0.45\% beyond \qty{10}{\micro\second}, and about 0.15\% beyond \qty{100}{\micro\second}.
\item Between \qty{10}{\micro\second} and \qty{100}{\micro\second}: 50.2\% of messages wait more than the first and 49.0\% more
than the second. About half wait nothing, and those that wait, wait milliseconds.
\item The wait grows with the square of a gap: a message is more likely to land in a long gap, and then waits longer.
\item $250 \times \qty{2}{\micro\second} = 0.05\%$ and $250 \times (\qty{2}{\micro\second})^2 / 2 = \qty{0.5}{\nano\second}$.
\item $250 \times (\qty{2}{\micro\second} - \qty{1}{\micro\second}) / \qty{1}{\second} = 0.025\%$.
\item $2 \times 10^{-6}$ of the time and \qty{0.002}{\nano\second} on average.
\item On a housekeeping CPU, which runs it on behalf of the tickless ones.
\item Up to \qty{50}{\milli\second} once a second, 5\% of the time: worse than the untuned laptop.
\item It is a virtual machine: its kernel has no adaptive ticks, the hypervisor's own work steals time the guest cannot see
or move, and an ordinary user cannot change the boot line.
\item \textbf{Named result.} Pinned alone on this laptop, a spinning thread is not running about 1.3\% of the time, with a worst
gap of about \qty{5.4}{\milli\second}; on the modelled tuned host, $2 \times 10^{-6}$ of the time, with a worst gap of
\qty{2}{\micro\second}: some six thousand times less time, a worst gap some 2\,700 times shorter.
\item \texttt{isolated}, \texttt{siblings}, \texttt{nohz\_full}, \texttt{rcu\_nocbs}, \texttt{irq\_affinity},
\texttt{cstate} and \texttt{memlock}: seven of eleven.
\item On an isolated spare CPU of the same host, configured like the hot ones, or on a hot CPU before the open; never beside a
hot thread.
\item Its cause: the meter sees that the CPU was taken, not by whom. Kernel tracing (interrupt counts per CPU in
\texttt{/proc/interrupts} before and after, the scheduler's trace events) attributes it.
\item The audit's report of the host, next to the placement plan, the kernel's version and its boot line.
\item It gives the hot path whole CPUs that the kernel leaves alone, at the price of cores, power and a host that must be
configured and checked for it.
\end{enumerate}
\end{solution}

\begin{solution}{iq:ll:linux-tuning:1}
\texttt{isolcpus} removes CPUs from load balancing so that only explicitly pinned threads run there; \texttt{nohz\_full} stops
the scheduling-clock tick on them while they run one task; \texttt{rcu\_nocbs} moves the kernel's deferred callbacks off them.
Each removes a different source of interruption, and the tick cannot stop while other tasks or callbacks still need the CPU.
\iqlookfor{three sources of noise and why each needs its own setting.}
\end{solution}

\begin{solution}{iq:ll:linux-tuning:2}
A regular period suggests the tick or a timer: check the kernel's tick rate and whether the CPU is tickless; compare interrupt
counts per CPU before and after a run; look for other runnable tasks on the CPU and for the kernel's per-CPU work; trace
scheduler events on that CPU. Confirm with a hiccup meter on the same CPU.
\iqlookfor{measure first, then attribute with counters and tracing.}
\end{solution}

\begin{solution}{iq:ll:linux-tuning:3}
Blocking costs nothing when idle but adds the wake-up, tens of microseconds on the chapter's machine. \omterm{def:ll:linux-tuning:polling}{Busy polling} costs a
spinning core and saves the wake-up (\qty{2}{\micro\second} against \qty{30}{\micro\second} here). \omterm{def:ll:linux-tuning:polling}{Kernel bypass} costs a
supported card, a separate network stack and a spinning core, and removes the kernel from the path. Choose by the \omterm{def:ll:where-latency-comes-from:budget}{latency
budget}, per thread.
\iqlookfor{costs on each side and the budget as the deciding rule.}
\end{solution}

\begin{solution}{iq:ll:linux-tuning:4}
Usually not, if they run alone on isolated cores: there is nothing to preempt. The risks are real-time throttling (a
\qty{50}{\milli\second} gap a second by default when another task is runnable), a runaway thread locking up a CPU, and
\omterm{def:ll:lock-free-programming:inversion}{priority inversion} with lower-priority threads it shares locks with (chapter~11).
\iqlookfor{throttling and starvation, not just ``higher priority is faster''.}
\end{solution}

\begin{solution}{iq:ll:linux-tuning:5}
On housekeeping CPUs of the card's socket: off the hot cores, so they are not interrupted, but on the card's node, so that the
kernel's work on the packet leaves it in the socket's shared cache near the threads that read it, not across the
interconnect.
\iqlookfor{isolation from the hot core, locality to the card.}
\end{solution}

\begin{solution}{iq:ll:linux-tuning:6}
The rare, long gaps that tuning did not remove: interrupts still allowed on the hot CPUs, a tick that is still running (a
second runnable task, a kernel without adaptive ticks), \omterm{def:ll:cpp-for-latency-i-memory:fault}{page faults} (memory not locked or pre-faulted), a hyper-threading
sibling in use, the hypervisor on a virtual machine. Measure the gaps on the hot CPU and audit the host.
\iqlookfor{the tail comes from rare events; the proposition's squares.}
\end{solution}
